\label{tab:strongscale}
\caption{\textbf{Strong scaling of the distributed GNN solver.}
Fixed 3D cavity, $\mathrm{Ra}=10^{6}$, $h_{128}$ ($\approx\!1.57\times10^{6}$ pressure DOFs
total), across one to eight GPUs; ``wall-clock speedup'' is the whole-SIMPLE-step speedup of
\Cref{sec:si-scaling-protocol}. The kernels strong-scale sub-linearly (block-Jacobi ms/PCG-it
$6.29\to3.09$, $2.04\times$ on eight GPUs) while block-Jacobi's pressure iterations grow
steeply with rank ($29.0\to170.6$, a factor $5.9$), so the whole-step speedup saturates,
peaking near $1.6\times$ at four GPUs and regressing at eight. The Nicolaides deflation cuts
iterations $29$--$35\%$ at $\approx\!10\%$ extra per-iteration cost and leads on iterations,
absolute wall and speedup at every multi-GPU count (eight-GPU wall $750$ against $829$~ms). It
turns over at eight as well ($1.92\times$ at four, $1.72\times$ at eight): the deflation
raises the curve but does not remove the saturation, which is set by iteration growth this
fixed $1.57\times10^{6}$-DOF problem cannot amortize past four ranks.}

\label{tab:weakscale}
\caption{\textbf{Weak scaling of the distributed GNN pressure solve.}
3D cavity, $\mathrm{Ra}=10^{6}$, $\approx\!1.5\times10^{6}$ pressure DOFs per GPU
(GPU-representative load; protocol of \Cref{sec:si-scaling-protocol}). ``bj eff'' is
block-Jacobi's weak efficiency on ms/PCG-it against one GPU (np1\,/\,npX); ``it/o'' is
it/outer. The one-degree-per-rank
Nicolaides deflation cuts pressure iterations $27$--$35\%$ (growing with GPU count) at
$\approx\!10\%$ extra per-iteration cost, a net pressure-solve speedup of $22$--$29\%$
(``Nic.\ net''). The per-GPU load follows the admissible cubic grids, so it is not constant
($1.31$--$1.57$M) and efficiency is reported both raw and per DOF (``bj eff/DOF''). The two
differ most at four GPUs, which carries $17\%$ less work per GPU than the baseline; per DOF
the curve is monotone. A separate
five-repetition sweep at the lighter $2\times10^{5}$-DOF/GPU load gives $0.65$
($4.5\!\to\!6.9$~ms/PCG-it), the fixed per-iteration communication being less amortized
there.}

\label{tab:weak-ceiling}
\caption{\textbf{Weak scaling of the two three-dimensional conjugate problems at the
single-GPU memory ceiling.} Two-level preconditioner, conjugate rotational-IPCS + RK3--CN,
$x$-slab. Weak efficiency is referenced to two GPUs, with the value against one GPU in the
last column (\Cref{sec:si-scaling-protocol}).
\emph{Heat sink:} fixed $\approx\!1.5\times10^{7}$ total DOF/GPU, the largest problem one H200
holds \emph{for this operator set}, the mesh growing with NPROC to $1.13\times10^{8}$ DOF on
eight GPUs; $\mathrm{Ra}=10^{6}$, 3000-step timing. The full-GPU load makes the one-to-two-GPU onset
mild ($1.32\times$, against $3.45\times$ at a light $2\times10^{6}$~DOF/GPU load). Its
eight-GPU weak efficiency at $1.1\times10^{8}$ DOF is $0.74$ against two GPUs and $0.56$
against one; the residual loss is the two-level global coarse correction growing with the
total problem.
\emph{Engine:} fixed $\approx\!6\times10^{7}$~DOF/GPU ($104$~GB, $73\%$ of one H200), grown
with NPROC; $\mathrm{Ra}=10^{5}$, median of three runs. It weak-scales comparably to the
fluid-filled heat sink; a run at a lighter $1.6\times10^{7}$~DOF/GPU load weak-scales far
worse (see text).}
